\begin{tabular}{>{\raggedright\arraybackslash}p{4.3cm}>{\raggedright\arraybackslash}p{10.4cm}}
\toprule
Campo & Descripción \\
\midrule
\texttt{pais} & Nombre del país \\
\texttt{iso3} & Código ISO 3166-1 alfa-3 \\
\texttt{anio} & Año de referencia \\
\texttt{indicador} & Nombre del indicador \\
\texttt{codigo\_indicador} & Código de la serie en la fuente original \\
\texttt{tipo} & Observado | Estimación de la fuente | Proyección de la fuente | Proyección propia \\
\texttt{modelo} & Especificación empleada; vacío salvo en proyección propia \\
\texttt{valor} & Valor del indicador \\
\texttt{unidad} & Unidad de medida \\
\texttt{cobertura\_institucional} & Perímetro institucional del dato \\
\texttt{fuente} & Fuente primaria \\
\texttt{nota} & Advertencia de construcción o comparabilidad, cuando aplica \\
\midrule
\textbf{ADVERTENCIA DE USO} & \small Las proyecciones de gasto en salud y pensiones (tipo = 'Proyección propia') se re-estimaron en el Producto 4 (Opción 1: ventana de validación 2016-2019, anterior al choque de 2020). Tras la corrección 2 serie(s) siguen excediendo el rango histórico y conservan advertencia textual. Las series con MASE > 2 respecto del método ingenuo (Chile: ingresos del gobierno general; Costa Rica: ingresos del gobierno general; Brasil: balance primario; Chile: balance primario; Costa Rica: balance primario; Panamá: balance primario; Costa Rica: deuda neta del gobierno general; Colombia: gasto público en salud; México: gasto público en salud; Brasil: gasto público en pensiones) se conservan en la base marcadas como no informativas y no deben citarse como estimación puntual. El balance primario no se grafica: es estacionario y su extrapolación univariada no aporta información más allá de su media. Las series 'Crecimiento nominal del PIB' y 'Balance fiscal global' se incorporaron del WEO abril 2025 para el cálculo del diferencial (r - g) de la Sección 6. \\
\bottomrule
\end{tabular}
